% Ad Atticum 8.3 --- headnote
% ad-atticum-08-03, 18 February 49 BC, Cales

\textit{Cicero to Atticus, written from Cales on the
twelfth day before the Kalends of March 49 BC (the
manuscript dateline: \textit{Scr. in Caleno xii K.
Mart. a. 705 (49)}). This is the dramatic centerpiece
of \textit{Ad Atticum} book 8 --- the long deliberation
letter in which Cicero, having decided in 7.26 that ``if
there is to be war, I have decided to be with Pompey,''
now actually tries the decision against the facts on
the ground. Pompey has fallen back to Apulia; Caesar
has Picenum; Domitius is bottled up at Corfinium and
calling for help. Within days Pompey will fail to
relieve him, Domitius will surrender, and the road to
Brundisium and across the Adriatic will be all that is
left.}

\textit{The letter is built as a formal
\textit{deliberatio} in two columns: section 2 sets out
the case for following Pompey out of Italy (gratitude,
friendship, the cause of the Republic, the
intolerableness of remaining under one man's power);
sections 3--4 set out the case against, and constitute
one of Cicero's most pitiless audits of Pompey's whole
political career --- the indictments lined up in
anaphoric \textit{ille \dots ille \dots ille}, from
nursing Caesar up in the first place to the present
``most disgraceful flight'' from Rome. Section 4 turns
to the practical impossibilities: the Lower Sea in
midwinter, his brother and son to think of, the
\textit{fasces laureati} of his still-unrelinquished
Cilician imperium as so many \textit{compedes}
(``chains'') he would have to drag aboard. Section 6
weighs the historical precedents for staying behind
under a tyrant --- L.~Philippus, L.~Flaccus, Q.~Mucius
Scaevola the Pontifex under Cinna, against Thrasybulus
the Athenian liberator --- and then circles back to the
embarrassment of those same fasces: Caesar might offer
a triumph, and to accept or to refuse would be equally
ruinous. Section 7 breaks the symmetry with a
newsflash: a despatch has just come in by night that
Caesar is at Corfinium, that Domitius has a firm army
and wants to fight, that Pompey cannot in honour leave
him in the lurch; there is even a faint hope from Spain
that Afranius is closing in on Trebonius in the
Pyrenees. The whole letter, Cicero says at the end, is
written \textit{sedatiore animo quam proxime
scripseram} --- ``with a steadier mind than I wrote
with last time'' --- and ends, characteristically, by
laying the whole question back in Atticus's lap. The
section numbering jumps from 4 to 6 (the editors leave
no §5) and the manuscript has a textual crux at the
sentence on accepting or refusing the offered
triumph.}
